\documentclass[12pt]{article}

\usepackage{amsmath}
\usepackage{amssymb}
\usepackage{graphicx}
\usepackage{hyperref}
\usepackage[margin=1in]{geometry}
\usepackage{natbib}
\usepackage{xcolor}
\usepackage{tcolorbox}
\usepackage{booktabs}
\usepackage{tabularx}
\usepackage{longtable}
\usepackage{lineno}
\linenumbers
\tcbuselibrary{breakable,skins}
\bibliographystyle{aasjournal}

\newcommand{\bk}{\mathbf{k}}
\newcommand{\bx}{\mathbf{x}}
\newcommand{\bv}{\mathbf{v}}
\newcommand{\Omm}{\Omega_m}
\newcommand{\lambdadB}{\lambda_{\rm dB}}
\newcommand{\mchi}{m_\chi}
\newcommand{\eV}{\,\mathrm{eV}}
\newcommand{\kpc}{\,\mathrm{kpc}}
\newcommand{\Mpc}{\,\mathrm{Mpc}}
\newcommand{\Msun}{\,M_\odot}
\newcommand{\kms}{\,\mathrm{km\,s^{-1}}}
\newcommand{\mtwtwo}{m_{22}}

\title{Fuzzy Dark Matter and BECDM}
\author{}
\date{\today}

\begin{document}

\maketitle
\tableofcontents
\newpage

\begin{tcolorbox}[colback=yellow!10!white, colframe=orange!75!black,
                  title={\bfseries Disclaimer}]
These notes were compiled by an AI assistant (Claude, Anthropic) and have
\textbf{not been reviewed by a human domain expert}.  They may contain
errors in derivations, physical reasoning, or numerical estimates.
Cross-check key equations against the primary literature before use.
\end{tcolorbox}
\vspace{0.5em}

%----------------------------------------------------------------------
\section{Introduction and motivation}
\label{sec:intro}

Standard cold dark matter (CDM) reproduces the CMB, large-scale
galaxy clustering, and cluster dynamics at the few-percent level.
Where it has been challenged is on sub-galactic scales:
missing-satellites, too-big-to-fail, and cusp-vs-core (see
\texttt{sidm.tex} for a detailed discussion).  One radical response is
to blame the DM microphysics --- perhaps dark matter is not a heavy
weakly-interacting particle but an \emph{ultralight bosonic field}
whose de Broglie wavelength is of order kpc so that quantum
uncertainty smooths out structure on those scales.  \citet{Hu2000}
introduced this as \emph{fuzzy cold dark matter} (FCDM/FDM).

Naming zoo (all the same physics):
\begin{itemize}
  \item \textbf{FDM/FCDM} --- \citet{Hu2000}, phenomenological
        astrophysics.
  \item \textbf{Ultralight axion(-like)} (ULA / ULDM) --- particle
        physics language, especially string-theory motivated
        \citep{Arvanitaki2010,Marsh2016}.
  \item \textbf{BECDM} --- condensed-matter language, emphasises the
        macroscopic-wavefunction / Bose--Einstein-condensate viewpoint
        \citep{Mocz2017,Mocz2020}.
  \item \textbf{SFDM / $\psi$DM / wave dark matter} --- generic
        agnostic names \citep{Schive2014,Schive2014b}.
\end{itemize}

The characteristic mass scale is $m\sim 10^{-22}\eV$, chosen so that
$\lambdadB\sim 1\,{\rm kpc}$ in a typical dwarf galaxy --- matching
the size of observed cores.  The comprehensive review is
\citet{Hui2017}; particle-physics side is \citet{Marsh2016}.

%----------------------------------------------------------------------
\section{Particle physics: ultralight scalars and misalignment}
\label{sec:particle}

\subsection{Scalar-field action and cosmological equation of motion}

Take a real scalar field $\phi$ minimally coupled to gravity:
%
\begin{equation}
  S = \int d^4x\,\sqrt{-g}\left[
       -\tfrac{1}{2}g^{\mu\nu}\partial_\mu\phi\,\partial_\nu\phi
       -\tfrac{1}{2}m^2\phi^2
      \right],
  \label{eq:scalar_action}
\end{equation}
%
with self-interactions omitted at leading order.  The
Euler--Lagrange equation is the Klein--Gordon equation
$\Box\phi = m^2\phi$.  In a flat FRW background with metric
determinant $\sqrt{-g} = a^3$, the covariant d'Alembertian gives
%
\begin{equation}
  \ddot\phi + 3H\dot\phi + m^2\phi = 0.
  \label{eq:KG_FRW}
\end{equation}

\subsection{Misalignment mechanism: cold-matter behaviour}

Two regimes of Eq.~\eqref{eq:KG_FRW}:

\paragraph{Overdamped ($H\gg m$).}  The friction term $3H\dot\phi$
dominates.  Any initial value $\phi_i$ is held approximately constant
(``frozen misalignment''), so $\rho_\phi\simeq V(\phi_i) = m^2\phi_i^2/2$
is constant and $w = -1$; the field acts as a small cosmological
constant.

\paragraph{Oscillating ($H\ll m$).}  Substitute the WKB ansatz
$\phi = A(t)\cos(mt+\theta)$ with $|\dot A/(mA)|\ll 1$.  To leading
order in $\dot A/(mA)$,
%
\begin{equation}
  \ddot\phi + 3H\dot\phi + m^2\phi
    \approx -2m\dot A\sin(mt+\theta) - 3HmA\sin(mt+\theta) = 0,
\end{equation}
%
which fixes $\dot A = -\tfrac{3}{2}HA$, i.e.\ $A\propto a^{-3/2}$.
The averaged energy density and pressure are
%
\begin{equation}
  \langle\rho_\phi\rangle = \tfrac{1}{2}m^2 A^2 \propto a^{-3},
  \qquad
  \langle p_\phi\rangle = 0
  \qquad \Longrightarrow\qquad w = 0,
\end{equation}
%
so the oscillating field redshifts as pressureless matter.

The transition occurs at $H(a_{\rm osc})\sim m/3$.  For $m\sim
10^{-22}\eV$ and a radiation-dominated era, this happens at
%
\begin{equation}
  T_{\rm osc} \sim (m\,M_{\rm Pl})^{1/2}
    \sim 500\eV,
  \qquad z_{\rm osc}\sim 10^{6},
\end{equation}
%
well before matter--radiation equality, so by $z=1000$ the field has
been oscillating for a long time and is indistinguishable from CDM on
super-de-Broglie scales.  Requiring today's abundance
$\Omega_\phi h^2\sim 0.12$ gives an ``axion decay constant''
$F_\phi\sim 10^{17}\,{\rm GeV}$ for the sweet-spot mass ---
comparable to the GUT/string scale.

\subsection{Where does $10^{-22}\eV$ come from?}

The QCD axion (original Peccei--Quinn resolution of the strong-CP
problem) has $m\sim 10^{-6}$--$10^{-3}\eV$, set by QCD instanton
effects and $f_a$.  \emph{Ultralight} axion-like particles arise
generically in string compactifications --- the string axiverse
\citep{Arvanitaki2010}: many moduli acquire tiny masses from
non-perturbative dynamics,
%
\begin{equation}
  m \sim M_{\rm Pl}\,e^{-S_{\rm inst}},
\end{equation}
%
with $S_{\rm inst}$ a large instanton action.  A small change in
$S_{\rm inst}$ shifts $m$ by many orders of magnitude, so a plausible
theoretical prior spans all of
$10^{-33}\eV$--$10^{-3}\eV$.  Astrophysically, $m\sim 10^{-22}\eV$ is
special because that is where $\lambdadB$ inside a dwarf galaxy
matches the observed core size.

\subsection{Bose occupation number}

At $m\sim 10^{-22}\eV$ and local dark-matter density
$\rho\sim 0.3\,{\rm GeV\,cm^{-3}}$, the number density is
$n = \rho/m\sim 10^{93}\,{\rm Mpc^{-3}}$, and the number of quanta per
de Broglie cell is $\sim 10^{96}$.  The field is highly degenerate and
Bose-condensed on any astrophysical scale, so it can be treated as a
classical complex scalar $\psi(\bx,t)$ whose $|\psi|^2$ is the DM
density.  ``BECDM'' is not a metaphor.

%----------------------------------------------------------------------
\section{Non-relativistic reduction: Schr\"odinger--Poisson}
\label{sec:sp}

\subsection{Derivation from Klein--Gordon}

Working in units $\hbar=c=1$ (restored below), consider the KG
equation on the perturbed FRW metric
$ds^2 = -(1+2\Phi)dt^2 + a^2(1-2\Phi)d\bx^2$
with Newtonian potential $|\Phi|\ll 1$.  Two scale hierarchies apply:
%
\begin{equation}
  |\dot\psi/\psi| \ll m,
  \qquad
  |\nabla\psi/\psi| \ll ma
  \qquad\text{(non-relativistic, sub-Compton).}
\end{equation}
%
Factor out the fast Compton oscillation with the ansatz
%
\begin{equation}
  \phi(\bx,t) = \frac{1}{\sqrt{2m}}
                \Bigl[\psi(\bx,t)\,e^{-imt} + \psi^*(\bx,t)\,e^{+imt}\Bigr],
  \label{eq:WKB_ansatz}
\end{equation}
%
promoting $\psi$ to a slowly varying complex envelope.  Substitute
\eqref{eq:WKB_ansatz} into $\Box\phi = m^2\phi$; the $m^2\phi$ terms
cancel the leading Compton-scale contribution, so what survives is
%
\begin{equation}
  2im\dot\psi + 3imH\psi + a^{-2}\nabla^2\psi - 2m^2\Phi\psi = 0,
\end{equation}
%
(neglecting $\ddot\psi$, $H^2\psi$, and $\Phi\dot\psi$).  Dividing by
$2m$ and restoring $\hbar$ gives the \emph{Schr\"odinger equation in
an expanding background}:
%
\begin{equation}
  i\hbar\Bigl(\partial_t + \tfrac{3}{2}H\Bigr)\psi
    = -\frac{\hbar^2}{2m\,a^2}\nabla^2\psi + m\Phi\psi.
  \label{eq:schrod_FRW}
\end{equation}
%
The Hubble drag $3H\psi/2$ is exactly what makes
$\rho = m|\psi|^2$ redshift as $a^{-3}$.  On timescales much shorter
than $H^{-1}$ we can drop $3H\psi/2$ and set $a=1$; the gravitational
potential is sourced by
%
\begin{equation}
  \nabla^2\Phi = 4\pi G\,\bigl(|\psi|^2 - \bar\rho\bigr).
  \label{eq:poisson_fdm}
\end{equation}
%
Equations \eqref{eq:schrod_FRW}--\eqref{eq:poisson_fdm} form the
\emph{Schr\"odinger--Poisson} (SP) system.  In natural units,
%
\begin{equation}
  i\hbar\,\partial_t\psi
    = -\frac{\hbar^2}{2m}\nabla^2\psi + m\Phi\psi,
  \qquad
  \nabla^2\Phi = 4\pi G\,(|\psi|^2-\bar\rho).
  \label{eq:SP}
\end{equation}
%
Reference: \citet{Hu2000}, \citet{Marsh2016}, \citet{Hui2017}.

\subsection{Madelung transform: hydrodynamic form}

Write $\psi$ in polar form,
%
\begin{equation}
  \psi(\bx,t) = \sqrt{\rho(\bx,t)/m}\;e^{iS(\bx,t)/\hbar},
  \label{eq:madelung_def}
\end{equation}
%
identifying the fluid velocity $\bv = \nabla S/m$.
Substituting into \eqref{eq:SP} and separating real and imaginary
parts yields two coupled equations.

\paragraph{Imaginary part (continuity).}  Multiplying the imaginary
part by $2\sqrt\rho/(\hbar m)$ gives
%
\begin{equation}
  \partial_t\rho + \nabla\!\cdot(\rho\bv) = 0.
  \label{eq:continuity_fdm}
\end{equation}

\paragraph{Real part (Bernoulli/Euler).}  The real part is
%
\begin{equation}
  -\partial_t S = \frac{(\nabla S)^2}{2m} + m\Phi
                  - \frac{\hbar^2}{2m}\frac{\nabla^2\sqrt\rho}{\sqrt\rho}.
\end{equation}
%
Take one spatial gradient and use $\bv=\nabla S/m$:
%
\begin{equation}
  \partial_t\bv + (\bv\cdot\nabla)\bv
    = -\nabla\Phi - \nabla Q,
  \label{eq:euler_fdm}
\end{equation}
%
with the \emph{quantum potential}
%
\begin{equation}
  Q(\bx,t) \equiv -\frac{\hbar^2}{2m^2}
                  \frac{\nabla^2\sqrt\rho}{\sqrt\rho}.
  \label{eq:Q_def}
\end{equation}

\paragraph{Physical content.}  Eqs.~\eqref{eq:continuity_fdm} and
\eqref{eq:euler_fdm} are the continuity and Euler equations of a
\emph{pressureless} self-gravitating fluid supplemented by a
non-local, gradient-based ``pressure'' $Q$.  $Q$ is not a genuine
thermodynamic pressure --- it depends on curvature of $\sqrt\rho$, not
on $\rho$ locally --- but it acts like one on scales below the de
Broglie length, opposing density gradients and preventing collapse.
Set $\hbar\to 0$ and one recovers pressureless CDM.  The FDM
``fuzziness'' is entirely $Q$.

Note that $\bv=\nabla S/m$ implies $\nabla\times\bv=0$ identically:
FDM velocity is irrotational.  Vorticity in FDM is
\emph{quantised} --- it lives on defect lines where $|\psi|=0$
(nodes/vortex tubes) where $S$ is multivalued.  This makes the
Madelung representation break at vortex lines, an issue relevant for
numerical implementation \citep{Mocz2017}.

%----------------------------------------------------------------------
\section{Linear theory: Jeans scale and quantum pressure}
\label{sec:jeans}

\subsection{Dispersion relation from Schr\"odinger--Poisson}

Linearise \eqref{eq:SP} around a homogeneous background of density
$\bar\rho$.  In a static (non-expanding) background, perturbations
$\delta\rho \propto e^{i\bk\cdot\bx + \gamma t}$ satisfy
%
\begin{equation}
  \gamma^2 = 4\pi G\bar\rho
             - \left(\frac{\hbar k^2}{2m}\right)^{\!2}.
  \label{eq:dispersion_fdm}
\end{equation}
%
The first term is standard Jeans gravity; the second, quartic in $k$,
is the quantum-pressure contribution inherited from
$-\hbar^2\nabla^2/(2m)$ acting on the amplitude \citep{Hu2000,Hui2017}.

\subsection{Jeans wavenumber}

Setting $\gamma=0$ in \eqref{eq:dispersion_fdm} gives
%
\begin{equation}
  k_J^4 = \frac{16\pi G m^2 \bar\rho}{\hbar^2}
  \quad\Longrightarrow\quad
  k_J = \left(\frac{16\pi G m^2\bar\rho}{\hbar^2}\right)^{\!1/4}.
  \label{eq:kJ_def}
\end{equation}
%
Modes with $k>k_J$ are quantum-pressure stabilised and oscillate;
modes with $k<k_J$ are gravitationally unstable and grow.  The
corresponding physical scale
%
\begin{equation}
  r_J = \frac{2\pi}{k_J}
      = \pi^{3/4}\,(G\bar\rho)^{-1/4}\,m^{-1/2},
\end{equation}
%
is the \emph{geometric mean} of the gravitational dynamical scale
$(G\bar\rho)^{-1/2}$ and the Compton scale $m^{-1}$.  In cosmological
units \citep{Hu2000},
%
\begin{equation}
  r_J = 55\,\mtwtwo^{-1/2}
          (\rho/\bar\rho_b)^{-1/4}
          (\Omm h^2)^{-1/4}\kpc,
  \qquad \mtwtwo \equiv m/10^{-22}\eV.
  \label{eq:rJ_kpc}
\end{equation}

\subsection{Uncertainty-principle interpretation}

An intuitive derivation: for a mode fitting inside a self-gravitating
region of density $\bar\rho$ and radius $r$, orbital velocity is
$v\sim(G\bar\rho)^{1/2}r$, so orbital momentum is $p\sim mv$.  The
de Broglie wavelength associated with $p$ is
$\lambda_{\rm dB}\sim\hbar/(mv)\sim\hbar\,(G\bar\rho)^{-1/2}(mr)^{-1}$.
Requiring $\lambda_{\rm dB}\sim r$ (mode ``fits'') gives
%
\begin{equation}
  r^2 \sim \frac{\hbar}{m}(G\bar\rho)^{-1/2}
  \;\Longrightarrow\;
  r \sim \hbar^{1/2}\,m^{-1/2}\,(G\bar\rho)^{-1/4},
\end{equation}
%
which is $r_J$ up to $O(1)$ factors.  So $r_J$ is the smallest scale
on which a wave-DM mode can be gravitationally bound; below this
scale Heisenberg uncertainty wins.

\subsection{Transfer function}

The cosmological transfer function is imprinted by the fact that the
comoving Jeans wavenumber $k_{J,{\rm com}}(a)$ evolves only weakly
during radiation and matter domination:
$k_{J,{\rm com}}\propto a^{1/4}$.  As a result the linear power
spectrum develops a \emph{sharp break} near $k_{J,{\rm eq}}$ (the
comoving Jeans scale at matter--radiation equality) rather than the
gradual $k^{-2}$ suppression of warm dark matter.  \citet{Hu2000} fit
%
\begin{equation}
  T_F(k) \simeq \frac{\cos(x^3)}{1+x^8},
  \qquad
  x = 1.61\,\mtwtwo^{1/18}\,\frac{k}{k_{J,{\rm eq}}},
  \label{eq:hbg_TF}
\end{equation}
%
with
%
\begin{equation}
  k_{J,{\rm eq}} \simeq 9\,\mtwtwo^{1/2}\,{\rm Mpc}^{-1}.
\end{equation}
%
Half-power wavenumber:
$k_{1/2}\simeq 4.5\,\mtwtwo^{4/9}\,{\rm Mpc}^{-1}$, at which the
transfer function drops by a factor of 2 in amplitude.

Once modes oscillate they never re-enter the growing branch, so the
suppression is baked in --- $T_F$ is effectively redshift-independent
below $k_{J,{\rm eq}}$ for the practical range $\mtwtwo\sim 1$
\citep{Schive2014b}.  This justifies the widely-used practice of
folding $T_F^2$ into the CDM $P(k)$ before running an $N$-body
simulation.

%----------------------------------------------------------------------
\section{Soliton cores}
\label{sec:soliton}

\subsection{Stationary Schr\"odinger--Poisson and scaling symmetry}

Look for stationary states of \eqref{eq:SP} of the form
$\psi(\bx,t) = \chi(\bx)\,e^{-i\gamma m t/\hbar}$.  Then $\chi$
satisfies the eigenvalue problem
%
\begin{equation}
  -\frac{\hbar^2}{2m}\nabla^2\chi + m\Phi\chi = m\gamma\chi,
  \qquad
  \nabla^2\Phi = 4\pi G m|\chi|^2.
  \label{eq:stationary_SP}
\end{equation}
%
The system admits a one-parameter scaling symmetry.  Under
$\lambda>0$,
%
\begin{equation}
  r\to\lambda^{-1}r,
  \quad
  \chi\to\lambda^2\chi,
  \quad
  \Phi\to\lambda^2\Phi,
  \quad
  \gamma\to\lambda^2\gamma,
  \quad
  M \equiv m\!\int|\chi|^2\to\lambda M,
  \label{eq:soliton_scaling}
\end{equation}
%
\eqref{eq:stationary_SP} is invariant.  In particular a bound-state
solution with mass $M$ has size $r\propto M^{-1}$, central density
$\rho_c\propto M^4$, and central potential
$|\Phi_c|\propto M^2$ \citep{Marsh2016,Hui2017}.

\subsection{Ground-state (soliton) profile}

The nodeless ground state of Eq.~\eqref{eq:stationary_SP} is called
the ``soliton'' (analogue of a boson star).  Its numerical eigenvalue
is $\gamma_1\simeq -0.69$ in scale-normalised units, and its density
profile is well approximated by \citep{Schive2014}
%
\begin{equation}
  \rho_{\rm sol}(r) \simeq
  \frac{\rho_c}{\bigl[1 + 0.091\,(r/r_c)^2\bigr]^8},
  \label{eq:soliton_fit}
\end{equation}
%
where $r_c$ is the half-density radius.  Reinstating dimensional
factors, the soliton core radius and central density obey
%
\begin{equation}
  r_c \simeq \frac{3.925\,\hbar^2}{G M m^2},
  \qquad
  \rho_c \propto \Bigl(\frac{Gm^2}{\hbar^2}\Bigr)^{\!3} M^4.
  \label{eq:soliton_scaling_dim}
\end{equation}
%
Equivalently: $r_c M \propto (\hbar/m)^2$, so heavier solitons are
smaller and denser --- opposite to CDM haloes.  A useful benchmark:
$r_c \sim 22 (\rho_c/\rho_{\rm crit})^{-1/4} \mtwtwo^{-1/2}\kpc$
\citep{Marsh2016}.

\subsection{Quantum virial theorem}

Compute the moment of inertia
$I \equiv \tfrac{1}{2}m\int d^3r\,r^2|\psi|^2$ of a stationary
Schr\"odinger--Poisson state.  Its second time derivative
$\ddot I$ can be evaluated using the SP Hamiltonian
$\hat H = -\hbar^2\nabla^2/(2m)+m\Phi$.  After manipulation
\citep{Hui2017},
%
\begin{equation}
  \ddot I = 2K + V + 2Q,
\end{equation}
%
where
%
\begin{align*}
  K &\equiv \frac{1}{2}\int d^3r\,\frac{\rho}{m}|\bv|^2
    &&\text{(bulk kinetic)}, \\
  V &\equiv -G\!\int d^3r\!\int d^3r'\,
    \frac{\rho(\bx)\rho(\bx')}{|\bx-\bx'|}
    &&\text{(gravitational PE)}, \\
  Q &\equiv \frac{\hbar^2}{2m^2}\!\int d^3r\,|\nabla\sqrt\rho|^2
    &&\text{(quantum kinetic energy)}.
\end{align*}
%
For a stationary state, $\ddot I = 0$, so $2K + V + 2Q = 0$.  The
soliton has $\bv=0$ everywhere (single-phase wavefunction), so $K=0$
and $V=-2Q$: the soliton saturates the quantum virial bound
$Q/|V| = 1/2$.

\subsection{Soliton--halo relation}

$N$-body/pseudo-spectral SP simulations \citep{Schive2014,Schive2014b}
show that virialised FDM haloes host a central soliton \emph{plus} an
NFW-like envelope.  \citet{Schive2014b} measured the empirical
relation
%
\begin{equation}
  M_{\rm sol} \simeq
    2.6\times 10^{9}\Msun\,
    \left(\frac{1+z}{1}\right)^{1/2}
    \left(\frac{M_{\rm halo}}{10^{12}\Msun}\right)^{1/3}
    \mtwtwo^{-1},
  \label{eq:M_soliton_halo}
\end{equation}
%
essentially: the soliton and its host halo share the same
specific binding energy, $|E|/M_{\rm sol} \simeq |E|/M_{\rm halo}$.
The scaling $M_{\rm sol}\propto M_{\rm halo}^{1/3}$ implies smaller
haloes host \emph{proportionally larger} soliton fractions --- 100\%
soliton in the tiniest dwarfs, negligible in clusters.

\citet{Bar2018} pointed out that this relation, combined with the
scaling \eqref{eq:soliton_scaling_dim}, forces a distinctive bump in
the rotation curve at the soliton radius: the ratio of the maximum
soliton circular velocity to the host halo's peak circular velocity is
$\sim 1$ independent of halo mass and $m$.  This is a strong
observational discriminant.

%----------------------------------------------------------------------
\section{Numerical methods}
\label{sec:numerical}

Two mainstream approaches to solving the SP system in cosmology.

\subsection{Pseudo-spectral (split-step) $\psi$-solver}

Split the time-evolution operator via Strang splitting:
%
\begin{equation}
  \psi(t+\Delta t)
    = e^{-iV\Delta t/(2\hbar)}\,
      e^{-iT\Delta t/\hbar}\,
      e^{-iV\Delta t/(2\hbar)}\,\psi(t),
\end{equation}
%
where $T = -\hbar^2\nabla^2/(2m)$ (kinetic) and $V = m\Phi$
(potential).  Apply $e^{-iT\Delta t/\hbar}$ in Fourier space (exact
plane-wave rotation) and $e^{-iV\Delta t/(2\hbar)}$ in real space
(pointwise multiplication).  Second-order accurate, unitary,
manifestly $L^2$-conserving.  \citep{Schive2014,MaySpringel2021,Mocz2017,Mocz2020}.

Time step constraints: two conditions must both be satisfied,
%
\begin{align}
  \Delta t &\lesssim \frac{2\pi\hbar}{m\,\max|\Phi|}
    &&\text{(phase advance, potential step)}, \\
  \Delta t &\lesssim \frac{4}{3\pi}\,\frac{m\,\Delta x^2}{\hbar}
    &&\text{(Nyquist kinetic step)}.
\end{align}
%
The kinetic constraint forces $\Delta t\propto \Delta x^2 \cdot m$
--- very restrictive at high spatial resolution and small $m$.

\citet{MaySpringel2021} implement this in AREPO on runs up to
$8640^3$ grid cells in a $10\Mpc/h$ box for $m=2\times 10^{-23}\eV$,
resolving the de Broglie scale in Milky Way-mass haloes.
Interference-driven granular substructure and central solitons are
recovered.

\subsection{Madelung / SPH-fluid ($\rho$-solver)}

Solve Eqs.~\eqref{eq:continuity_fdm}--\eqref{eq:euler_fdm} treating
$-\nabla Q$ as an extra body force in an SPH or moving-mesh
hydrodynamic solver.  \citet{Mocz2017} implement this in AREPO.
Advantages: reuse existing hydro machinery, coarser grids acceptable
(only $\nabla^2\sqrt\rho$ needs resolution).  Disadvantage: the
Madelung representation breaks at vortex lines where $\rho\to 0$;
these must be treated specially with a hybrid ``switch to $\psi$''
prescription \citep{Mocz2020}.

\citet{NoriBaldi2018} take an intermediate path: modify GADGET-3's
SPH to include a fully consistent quantum-pressure tensor derived
from the SP moments, yielding AX-GADGET, with cosmological scaling to
$\sim 512^3$ particles for $m\gtrsim 10^{-22}\eV$.

\subsection{Practical scaling}

Simulation cost scales as $t_{\rm sim}\propto \Delta x^{-4}$ (spatial
plus temporal) --- prohibitive at high resolution.  For present-day
compute, the practical range is
$m\gtrsim 3\times 10^{-23}\eV$ in $\sim 10\Mpc/h$ boxes at kpc-scale
resolution.  Massive parallel-GPU pseudo-spectral (GAMER
\citep{Schive2014}, AREPO \citep{MaySpringel2021}) or hybrid Madelung
codes are pushing the frontier.

%----------------------------------------------------------------------
\section{Astrophysical signatures}
\label{sec:signatures}

\subsection{Halo mass function cutoff}

The transfer-function cutoff at $k_{J,{\rm eq}}$
(Eq.~\ref{eq:hbg_TF}) suppresses the halo mass function below the
half-mode mass
%
\begin{equation}
  M_{1/2} \simeq 3.8\times 10^{10}\,\mtwtwo^{-4/3}\Msun.
\end{equation}
%
For $\mtwtwo\sim 1$ this is close to the missing-satellite scale ---
motivating $10^{-22}\eV$ as ``the fix''.

\subsection{Soliton cores}

Every FDM halo hosts a central soliton whose mass follows
Eq.~\eqref{eq:M_soliton_halo}.  For a dwarf spheroidal
($M_{\rm halo}\sim 10^9\Msun$, $\mtwtwo\sim 1$), the soliton has
$M_{\rm sol}\sim 3\times 10^7\Msun$ and $r_c\sim 200\,{\rm pc}$ ---
converting a CDM cusp into a flat core without invoking baryonic
feedback \citep{Schive2014}.

\subsection{Granular substructure and dynamical heating}

Outside the soliton, the halo is a random superposition of many
excited modes.  Interference produces order-unity density
fluctuations on the de Broglie scale, with coherence time
$\tau_{\rm dB}\sim(mv^2)^{-1}\hbar$.  These granules act as
``quasiparticles'' of effective mass
$m_{\rm eff}\sim\rho(\lambdadB/2)^3$, driving two-body relaxation of
embedded objects (star clusters, subhaloes) with timescale
\citep{Hui2017}
%
\begin{equation}
  t_{\rm relax}
   \sim 10^{10}\,{\rm yr}\,
     \left(\frac{v}{100\kms}\right)^{\!2}
     \left(\frac{r}{5\kpc}\right)^{\!4}
     \mtwtwo^{3}\,f_{\rm relax}^{-1}.
\end{equation}
%
The same relaxation drives \emph{gravitational cooling}: kinetic
energy from the halo envelope is transferred inward, growing the
central soliton at the expense of the outer halo
\citep{Mocz2017,Mocz2020}.

\subsection{Tunneling floor on subhalo survival}

The soliton wavefunction has exponentially decaying tails at any
finite energy, so a subhalo of soliton mass $M$ in a host of mean
density $\bar\rho_{\rm host}$ can tunnel outside its tidal radius over
time.  Requiring survival for $\gtrsim 10$ orbits sets a minimum
subhalo central density $\rho_c\gtrsim 60\bar\rho_{\rm host}$, hence a
minimum surviving mass \citep{Hui2017}
%
\begin{equation}
  M \gtrsim 6.7\times 10^{8}\Msun\,
    \left(\frac{\mathcal{M}}{10^{11}\Msun}\right)^{1/4}
    \left(\frac{10\kpc}{a}\right)^{3/4}
    \mtwtwo^{-3/2}.
\end{equation}

%----------------------------------------------------------------------
\section{Observational constraints}
\label{sec:constraints}

Bounds on the ULA mass, ordered by tightening decade.  For the
``canonical FDM'' interpretation --- all of dark matter is a single
ultralight species --- the current bounds are:

\begin{tabularx}{\linewidth}{@{}l l X@{}}
\toprule
Probe & Bound $m$ [eV] & Reference \\
\midrule
CMB linear (Planck)         & $\gtrsim 10^{-24}$
   & \citet{Marsh2016} (weak) \\
Cosmic dawn / 21-cm         & $\gtrsim 10^{-22}$
   & (model-dependent) \\
UV luminosity function $z\!\sim\!4$--$10$
                            & $\gtrsim 1.2\times 10^{-22}$
   & \citet{Schive2014b} (2$\sigma$) \\
Milky Way subhalo counts    & $\gtrsim 3\times 10^{-21}$
   & (astrophysical) \\
Lyman-$\alpha$ forest       & $\gtrsim 2\times 10^{-20}$
   & \citet{Rogers2021} \\
Ultra-faint dwarf sizes + kinematics
                            & $\gtrsim 3\times 10^{-19}$
   & \citet{Dalal2022} \\
\bottomrule
\end{tabularx}

\paragraph{Reading the table.}  Ly$\alpha$ (\citealt{Rogers2021})
rules $m<2\times 10^{-20}\eV$ from small-scale $P(k)$ suppression;
UFD kinematics (\citealt{Dalal2022}) push $m$ up by another
$1.5$ dex from the required soliton-core sizes in Segue 1, Draco,
etc.  Both, taken at face value, exclude the ``sweet spot''
$m\sim 10^{-22}\eV$ by 2--3 orders of magnitude.

\paragraph{Escape hatches.}  (i) FDM as a \emph{subdominant} DM
fraction $f_{\rm FDM}\lesssim 0.1$ --- the dominant CDM continues to
form structure normally, FDM only supplies large-scale features.
(ii) A distribution of masses (string axiverse) rather than a single
$m$: the cutoff is smeared and the strongest bound applies only to
the fraction at that mass.  (iii) Attractive self-interactions
($\lambda\phi^4$ with $\lambda>0$) modify the soliton profile,
altering the UFD prediction \citep{Mocz2023}.

%----------------------------------------------------------------------
\section{Summary and open directions}
\label{sec:summary}

Fuzzy dark matter is the same particle-physics idea (ultralight
bosonic field populated by misalignment) viewed either as
(i) a linear-theory cosmological cutoff at
$k_{J,{\rm eq}}\sim 9\,\mtwtwo^{1/2}\,{\rm Mpc}^{-1}$, or (ii) a
non-linear condensate obeying Schr\"odinger--Poisson with
quantum-pressure cores and interference substructure.  Both pictures
are the same equations, viewed under two different microscopes.

The core mass scale $m\sim 10^{-22}\eV$ was chosen historically to
resolve small-scale CDM problems.  Ly$\alpha$ and dwarf-galaxy
kinematics have since pushed the allowed mass up by $\sim 3$ decades,
so the original ``fuzzy sweet spot'' is now viable only as a
subdominant fraction or with attractive self-interactions.  Where FDM
remains most alive:
%
\begin{itemize}
  \item As a probe of the string axiverse at masses too small for
        laboratory searches.
  \item As a $\lesssim 10\%$ ultralight component with distinctive
        interference signatures on small scales.
  \item As a numerical framework (Schr\"odinger--Poisson) for
        regularising cold-matter simulations regardless of the
        underlying microphysics.
\end{itemize}

%======================================================================
\appendix

\section{Reading the scalar-field action: a QFT primer}
\label{app:qft_primer}

The starting point of \S\ref{sec:particle} is Eq.~\eqref{eq:scalar_action}:
%
\begin{equation}
  S = \int d^4x\,\sqrt{-g}\left[
       -\tfrac{1}{2}\,g^{\mu\nu}\,\partial_\mu\phi\,\partial_\nu\phi
       -\tfrac{1}{2}\,m^2\phi^2
      \right].
  \tag{Eq.~\ref{eq:scalar_action}}
\end{equation}
%
This appendix unpacks every symbol for readers from adjacent fields
where scalar-field actions are not routine.  Nothing below is
FDM-specific --- it is the standard QFT boilerplate that underlies
all the physics of a fundamental scalar field.  An expanded,
self-contained treatment --- including Noether's theorem, canonical
quantisation, coherent states, and the classical-wave limit that
justifies treating FDM as a classical field --- is in
\texttt{qft.tex}.

\subsection{The principle of least action}
\label{app:least_action}

\paragraph{Mechanics.}  Given a Lagrangian $L(q,\dot q,t)$ for a
system with a single generalised coordinate $q(t)$, the action is
%
\begin{equation}
  S[q] = \int_{t_1}^{t_2}\! L(q(t),\dot q(t), t)\,dt.
  \label{eq:action_mech}
\end{equation}
%
The physical trajectory is the one that makes $S$ stationary
under variations $q(t)\to q(t) + \delta q(t)$ with fixed endpoints
$\delta q(t_1) = \delta q(t_2) = 0$.

\emph{Derivation of Euler--Lagrange in one line.}  Compute
$\delta S$ by chain rule:
%
\begin{equation}
  \delta S
    = \int_{t_1}^{t_2}\!
        \left[
          \frac{\partial L}{\partial q}\,\delta q
          + \frac{\partial L}{\partial\dot q}\,\delta\dot q
        \right] dt.
\end{equation}
%
Use $\delta\dot q = \frac{d}{dt}\delta q$ (variations and time
derivative commute).  Integrate the second term by parts:
%
\begin{equation}
  \int_{t_1}^{t_2}\!\frac{\partial L}{\partial\dot q}\,
    \frac{d(\delta q)}{dt}\,dt
   = \underbrace{\left[\frac{\partial L}{\partial\dot q}\,\delta q\right]_{t_1}^{t_2}}_{=0\;\text{by BCs}}
     - \int_{t_1}^{t_2}\!\frac{d}{dt}\!\left(\frac{\partial L}{\partial\dot q}\right)\,\delta q\,dt.
\end{equation}
%
The boundary term drops.  So
%
\begin{equation}
  \delta S = \int_{t_1}^{t_2}\!
     \left[
       \frac{\partial L}{\partial q}
       - \frac{d}{dt}\!\frac{\partial L}{\partial\dot q}
     \right]\delta q\,dt.
\end{equation}
%
Requiring $\delta S = 0$ for arbitrary $\delta q(t)$ forces the
bracketed quantity to vanish pointwise (variational lemma),
giving the \textbf{Euler--Lagrange equation}
%
\begin{equation}
  \frac{d}{dt}\!\frac{\partial L}{\partial\dot q}
    - \frac{\partial L}{\partial q} = 0.
  \label{eq:EL_mech}
\end{equation}
%
Sanity check: $L = \tfrac{1}{2}m\dot q^2 - V(q)$ gives
$\partial L/\partial\dot q = m\dot q$, $\partial L/\partial q =
-V'(q)$, so \eqref{eq:EL_mech} becomes $m\ddot q = -V'(q)$ --- Newton's
second law.

\paragraph{Generalisation to fields.}  Replace the finite set of
generalised coordinates $\{q_i(t)\}$ by a field
$\phi(\bx, t)$ defined at every spacetime point.  Instead of one
value per time, the ``coordinate'' is an infinite-dimensional
configuration --- one value per spatial point.  Split the Lagrangian
$L$ into a spatial integral of a local \emph{Lagrangian density}
$\mathcal L$:
%
\begin{equation}
  L(t) = \int d^3\bx\,\mathcal L\bigl(\phi(\bx,t),\,
                                        \partial_\mu\phi(\bx,t)\bigr),
  \qquad
  S[\phi]
    = \int dt\,L(t)
    = \int d^4x\,\mathcal L.
  \label{eq:action_field}
\end{equation}
%
The density $\mathcal L$ depends on $\phi$ and its first derivatives
(both time and space), lumped into $\partial_\mu\phi$ for brevity.
$d^4x = dt\,d^3\bx$ is the flat-space spacetime volume element; on a
curved background one inserts $\sqrt{-g}$ (\S\ref{app:metric}).

\emph{Derivation of field-theory Euler--Lagrange.}  Under a variation
$\phi(x)\to\phi(x)+\delta\phi(x)$ with $\delta\phi$ vanishing on the
spacetime boundary,
%
\begin{equation}
  \delta S
    = \int d^4x\,\left[
      \frac{\partial\mathcal L}{\partial\phi}\,\delta\phi
      + \frac{\partial\mathcal L}{\partial(\partial_\mu\phi)}\,
        \delta(\partial_\mu\phi)
    \right].
\end{equation}
%
Because variation and derivative commute, $\delta(\partial_\mu\phi) =
\partial_\mu(\delta\phi)$.  Integrate the second term by parts (4D
divergence theorem):
%
\begin{equation}
  \int d^4x\,\frac{\partial\mathcal L}{\partial(\partial_\mu\phi)}\,
    \partial_\mu(\delta\phi)
   = \underbrace{\int_{\partial\Omega}\!
     dS_\mu\,\frac{\partial\mathcal L}{\partial(\partial_\mu\phi)}\,\delta\phi}_{=0\;\text{by BCs}}
     - \int d^4x\,\partial_\mu\!
       \left[\frac{\partial\mathcal L}{\partial(\partial_\mu\phi)}\right]\,\delta\phi.
\end{equation}
%
So
%
\begin{equation}
  \delta S = \int d^4x\,
     \left[
       \frac{\partial\mathcal L}{\partial\phi}
       - \partial_\mu\!\frac{\partial\mathcal L}{\partial(\partial_\mu\phi)}
     \right]\delta\phi.
\end{equation}
%
Requiring $\delta S = 0$ for arbitrary $\delta\phi(x)$ gives the
\textbf{field-theory Euler--Lagrange equation}
%
\begin{equation}
  \partial_\mu\!\frac{\partial\mathcal L}{\partial(\partial_\mu\phi)}
    - \frac{\partial\mathcal L}{\partial\phi} = 0.
  \label{eq:field_EL}
\end{equation}
%
This is directly analogous to \eqref{eq:EL_mech}, but with a
spacetime divergence $\partial_\mu(\cdots)$ instead of the mechanics
time derivative $\frac{d}{dt}(\cdots)$, and with functional
derivative $\partial\mathcal L/\partial\phi$ instead of ordinary
$\partial L/\partial q$.

\emph{Sanity check: free scalar.}  With
$\mathcal L = -\tfrac{1}{2}\eta^{\mu\nu}\partial_\mu\phi\partial_\nu\phi
- \tfrac{1}{2}m^2\phi^2$ (Minkowski signature $(-,+,+,+)$),
%
\begin{align}
  \frac{\partial\mathcal L}{\partial\phi} &= -m^2\phi, \\
  \frac{\partial\mathcal L}{\partial(\partial_\mu\phi)}
    &= -\eta^{\mu\nu}\partial_\nu\phi
    = -\partial^\mu\phi.
\end{align}
%
Substituting into \eqref{eq:field_EL}:
%
\begin{equation}
  \partial_\mu(-\partial^\mu\phi) - (-m^2\phi) = 0
  \;\Longrightarrow\;
  -\Box\phi + m^2\phi = 0
  \;\Longrightarrow\;
  \Box\phi = m^2\phi,
\end{equation}
%
the Klein--Gordon equation.  On a curved metric, the analog
manipulation (using $\sqrt{-g}$) gives the covariant Klein--Gordon
equation \eqref{eq:KG_covariant}.

\subsection{Index notation and Einstein summation}
\label{app:index}

In relativistic physics, spacetime points have four coordinates:
$x^\mu = (t, \mathbf x) = (x^0, x^1, x^2, x^3)$, with $\mu\in\{0,1,2,3\}$
running over time (index 0) and three spatial directions.  Repeated
indices are automatically summed --- \emph{Einstein summation
convention}:
%
\begin{equation}
  A^\mu B_\mu \equiv \sum_{\mu=0}^{3}\,A^\mu B_\mu.
\end{equation}
%
Two flavours of index:
%
\begin{itemize}
  \item \textbf{Contravariant} $A^\mu$ (index up): components of a
        vector, e.g.\ $x^\mu$, $p^\mu$.
  \item \textbf{Covariant} $A_\mu$ (index down): components of a
        covector (1-form).
\end{itemize}
%
The metric tensor $g_{\mu\nu}$ raises and lowers indices:
$A_\mu = g_{\mu\nu} A^\nu$, $A^\mu = g^{\mu\nu} A_\nu$, with
$g^{\mu\nu}$ the inverse metric ($g^{\mu\alpha}g_{\alpha\nu} =
\delta^\mu_\nu$).  Contracted upper--lower pairs like
$g^{\mu\nu}\partial_\mu\phi\partial_\nu\phi$ are Lorentz scalars ---
invariant under coordinate transformations.

The partial derivative operator is
%
\begin{equation}
  \partial_\mu \equiv \frac{\partial}{\partial x^\mu}
    = \left(\frac{\partial}{\partial t},\,\nabla\right),
\end{equation}
%
and $\partial^\mu = g^{\mu\nu}\partial_\nu$ raises its index.

\subsection{The metric $g_{\mu\nu}$ and $\sqrt{-g}$}
\label{app:metric}

The metric tensor $g_{\mu\nu}(x)$ encodes the geometry of spacetime.
It defines the invariant interval between neighbouring points:
%
\begin{equation}
  ds^2 = g_{\mu\nu}\,dx^\mu\,dx^\nu.
\end{equation}
%
In flat (Minkowski) spacetime with signature $(-,+,+,+)$:
%
\begin{equation}
  g_{\mu\nu} = \eta_{\mu\nu} = \mathrm{diag}(-1, +1, +1, +1),
  \qquad
  ds^2 = -dt^2 + d\mathbf x^2.
\end{equation}
%
On the flat FRW background used in cosmology,
$g_{\mu\nu} = \mathrm{diag}(-1, a^2, a^2, a^2)$ with $a(t)$ the scale
factor.  Comoving spatial derivatives inherit a factor $a^{-2}$ when
converted to physical derivatives.

The quantity $\sqrt{-g} \equiv \sqrt{-\det(g_{\mu\nu})}$ is the
\emph{invariant volume element} Jacobian.  In Minkowski
$\sqrt{-g} = 1$; in FRW $\sqrt{-g} = a^3$.  Its role in the action is
to make $d^4x\,\sqrt{-g}$ a coordinate-invariant volume element ---
if you change coordinates, $\sqrt{-g}$ transforms to compensate for
the change in $d^4x$, keeping the physical volume the same.

\subsection{Scalar field $\phi$ and Lorentz invariance}
\label{app:scalar}

A \emph{scalar field} $\phi(x)$ takes a single real value at every
spacetime point.  Under Lorentz transformations $x\to x'$, a scalar
transforms trivially: $\phi'(x') = \phi(x)$.  This is the simplest
possible relativistic field (contrast: vector field, spinor field,
tensor field).

The gradient $\partial_\mu\phi$ is a covector (1-form).  Its
Lorentz-invariant square is
%
\begin{equation}
  g^{\mu\nu}\,\partial_\mu\phi\,\partial_\nu\phi
    = -(\partial_t\phi)^2 + (\nabla\phi)^2/a^2
  \qquad \text{(FRW, signature $-,+,+,+$)}.
\end{equation}
%
This is the \emph{kinetic term} for the field --- it plays the role
of $\dot q^2$ in mechanics, but relativistically it mixes time and
space derivatives with the appropriate signs.

\subsection{Anatomy of Eq.~\eqref{eq:scalar_action}}
\label{app:action_anatomy}

Now put the pieces together:
%
\begin{equation*}
  \underbrace{\int d^4x}_{\text{spacetime integral}}
  \;
  \underbrace{\sqrt{-g}}_{\text{invariant Jacobian}}
  \;\Bigl[
  \underbrace{-\tfrac{1}{2}g^{\mu\nu}\partial_\mu\phi\,\partial_\nu\phi}_{\text{kinetic term}}
  \;
  \underbrace{-\tfrac{1}{2}m^2\phi^2}_{\text{mass / potential term}}
  \Bigr].
\end{equation*}
%
Reading it in words:
%
\begin{itemize}
  \item $\int d^4x$: sum contributions over all spacetime.
  \item $\sqrt{-g}$: absorb coordinate-choice-dependent volume factors
        so the result is a coordinate-invariant number.
  \item $-\tfrac{1}{2}g^{\mu\nu}\partial_\mu\phi\,\partial_\nu\phi$:
        kinetic energy density.  Positive when $\phi$ varies with
        time or space; the overall minus sign combined with the
        $g^{00}=-1$ signature makes the time-derivative term have
        the right sign to be a canonical kinetic energy
        $\tfrac{1}{2}\dot\phi^2$ (after the sign contract with
        $g^{00}$).
  \item $-\tfrac{1}{2}m^2\phi^2$: mass / potential term.  This is
        the harmonic-oscillator potential $V(\phi) = \tfrac{1}{2}m^2\phi^2$
        that gives the field a rest-mass energy.
\end{itemize}

Applying the field-theory Euler--Lagrange equation \eqref{eq:field_EL}
to this Lagrangian density on a general metric gives the covariant
Klein--Gordon equation
%
\begin{equation}
  \Box\phi \equiv \frac{1}{\sqrt{-g}}\partial_\mu\!\bigl(\sqrt{-g}\,g^{\mu\nu}\partial_\nu\phi\bigr)
    = m^2\phi.
  \label{eq:KG_covariant}
\end{equation}
%
On flat spacetime this reduces to $(\Box\phi) = -\partial_t^2\phi +
\nabla^2\phi = m^2\phi$; on FRW with $\sqrt{-g}=a^3$ and $g^{00}=-1$,
$g^{ii}=a^{-2}$, one gets
%
\begin{equation}
  \ddot\phi + 3H\dot\phi - \frac{1}{a^2}\nabla^2\phi + m^2\phi = 0.
\end{equation}
%
Dropping the spatial-gradient term (for a homogeneous background)
reproduces Eq.~\eqref{eq:KG_FRW}.

\subsection{Why the sign and factor conventions}
\label{app:conventions}

Two conventions worth flagging because they can trip up readers
importing formulas from different references:
%
\begin{itemize}
  \item \textbf{Metric signature.}  ``Mostly plus'' $(-,+,+,+)$ is
        used by most GR and cosmology references (Wald, Weinberg,
        MTW).  Particle physics often uses $(+,-,-,-)$ (Peskin \&
        Schroeder).  The kinetic term flips sign between the two,
        so you'll see either $-\tfrac{1}{2}g^{\mu\nu}\partial_\mu\phi\partial_\nu\phi$
        or $+\tfrac{1}{2}g^{\mu\nu}\partial_\mu\phi\partial_\nu\phi$ ---
        both mean the same thing in the respective convention.
  \item \textbf{Factor of 2 in kinetic and mass terms.}
        The $\tfrac{1}{2}$ prefactors are conventional and get you
        canonical normalisation of $\phi$ in the equations of motion.
        Non-canonically normalised fields differ by a rescaling
        $\phi\to c\phi$ absorbed into effective mass or coupling.
\end{itemize}

\subsection{Minimally-coupled vs.\ non-minimally-coupled}
\label{app:minimal}

Equation \eqref{eq:scalar_action} is the action for a
\emph{minimally coupled} scalar field: it interacts with gravity
only through $\sqrt{-g}$ and $g^{\mu\nu}$, i.e.\ through the metric
in the standard combinations.  A \emph{non-minimally-coupled} scalar
would add a term like $\tfrac{1}{2}\xi R\phi^2$, where $R$ is the
Ricci scalar (a curvature invariant).  FDM as usually formulated is
minimally coupled --- the field just senses the geometry, without
directly coupling to curvature.  Non-minimally-coupled scalars
appear in inflation (Higgs inflation, generalised Galileons) but
not in the FDM literature.

\subsection{Summary}
\label{app:qft_summary}

Every piece of Eq.~\eqref{eq:scalar_action} has a mechanical analogue:
%
\begin{center}
\begin{tabular}{@{}ll@{}}
\toprule
Mechanics                          & Field theory \\
\midrule
Coordinate $q(t)$                  & Field $\phi(\mathbf x, t)$ \\
Velocity $\dot q$                  & Gradient $\partial_\mu\phi$ \\
Kinetic energy $\tfrac{1}{2}m\dot q^2$ & Kinetic density
                                     $-\tfrac{1}{2}g^{\mu\nu}\partial_\mu\phi\partial_\nu\phi$ \\
Potential $V(q)$                   & Potential density $V(\phi)=\tfrac{1}{2}m^2\phi^2$ \\
Lagrangian $L = T - V$             & Lagrangian density $\mathcal L$ \\
Action $S=\int L\,dt$              & Action $S = \int d^4x\,\sqrt{-g}\,\mathcal L$ \\
Newton's law from $\delta S=0$     & Klein--Gordon from $\delta S=0$ \\
\bottomrule
\end{tabular}
\end{center}

Everything else in \S\ref{sec:particle} follows from applying this
action variationally on the FRW background, then taking the
non-relativistic limit.

\bibliography{references}

\end{document}
